\par\noindent\textbf{(1)} 令 ρ=ρ\_c θ\textasciicircum{}n、r=αξ，α²=(n+1)Kρ\_c\textasciicircum{}\{1/n−1\}/(4πG)。用静水方程消去质量，得到 (1/ξ²)d/dξ(ξ²θ')=−θ\textasciicircum{}n，正则中心条件 θ(0)=1、θ'(0)=0，避免原点的表观奇性。
\par\noindent\textbf{(2)} 设 θ=1+aξ²+bξ⁴+O(ξ⁶)，左边为 6a+20bξ²+…，右边为 −[1+naξ²+…]。逐阶比较得 a=−1/6、b=n/120，所以 θ=1−ξ²/6+nξ⁴/120+O(ξ⁶)。
\par\noindent\textbf{(3)} 若第一个零点 ξ\_1 有限，恒星半径 R=αξ\_1，质量由积分及方程得 M=4πα³ρ\_c[−ξ\_1²θ'(ξ\_1)]。所以 R∝K\textasciicircum{}\{1/2\}ρ\_c\textasciicircum{}\{(1−n)/(2n)\}，M∝K\textasciicircum{}\{3/2\}ρ\_c\textasciicircum{}\{(3−n)/(2n)\}，括号内是只依赖 n 的结构常数。
\par\noindent\textbf{(4)} 当 n=3，质量指数 (3−n)/(2n)=0，即固定 K 的多方星质量与 ρ\_c 无关，体现相对论性简并压强极限的临界尺度。n=5 的解析 θ=(1+ξ²/3)\textasciicircum{}\{−1/2\} 永不在有限半径变零，故有限质量但半径无穷；不能把有限 ξ\_1 公式用于 n≥5。
